\section{Positive constraints and truncation near a finite set}
\label{sec:quasilocal}

The global gap of $H$ will be used through positive operators that annihilate
its ground vector separately, retain a gap when summed, and have rapidly
decaying graph-local tails.  These exact constraints supply the contraction
products in Section~\ref{sec:amplification}.  For the transport and scanner
arguments in Sections~\ref{sec:transport} and~\ref{sec:scanner}, we also
truncate near a prescribed finite set $S_0$.  The terms crossing cuts inside
$S_0$ then have controlled graph-ball supports, while the total truncation
error depends on $|S_0|$ rather than the surrounding volume.  Both
constructions use only the gap of the full Hamiltonian.

Write $H=\sum_{i\in\mathcal I}h_i$, where $i$ indexes the distinct nonempty
designated supports $X_i$ from Section~\ref{sec:prelim}, and choose an anchor
$a_i\in X_i$.  Thus $\diam_{d_\Lambda}X_i\le R$ and $\|h_i\|\le J$.
We use $N_l(a)=\{x\in\Lambda:d_\Lambda(a,x)\le l\}$ for the graph ball,
with integer $l\ge0$.  The normalized conditional expectation onto this ball
is the map
\begin{equation}
 E_{i,l}(B)=q^{-|\Lambda\setminus N_l(a_i)|}
       \Tr_{\Lambda\setminus N_l(a_i)}(B)
       \otimes I_{\Lambda\setminus N_l(a_i)}.
 \label{eq:quasilocal-ce}
\end{equation}
Tensor factors in this formula are placed at their original sites.
Equivalently, $E_{i,l}$ averages conjugation by independent on-site unitaries
outside the ball.  It is positive, unital, contractive in operator norm, and
fixes every operator supported in the ball.

The counting bounds from Section~\ref{sec:prelim} give at most $v_R$
sites in each original support, at most $\mu_R$ supports containing a
specified site, and at most $\mu_R$ labels with any specified anchor.
In particular,
\begin{equation}
 \sup_{x\in\Lambda}\sum_{i:x\in X_i}\|h_i\|\le b_0:=J\mu_R,
 \qquad |X_i|\le v_R.
 \label{eq:quasilocal-budget}
\end{equation}
All constants in this section are independent of $\Lambda$; they may depend
on $q,R,J,\Delta$, and on explicitly fixed exponents or size bounds.

\subsection{Propagation in graph distance and a smooth spectral filter}

We first record the dynamical estimate needed to localize a spectral filter.
This also makes clear why holes and separate connected components cause no
change in the constants.  Starting from the Lieb--Robinson
propagation framework~\cite{LiebRobinson1972}, we use the finite-volume
argument of Nachtergaele, Ogata, and Sims~\cite{NachtergaeleOgataSims2006}, giving the
commutator iteration here in the graph metric of $\Lambda$.

\begin{lemma}[Graph-distance propagation]
\label{lem:quasilocal-lr}
Let $\tau_t(B)=e^{itH}Be^{-itH}$.  There are positive constants $C,v,c$
such that, for every original support $X_i$, every operator $A$ supported
in $X_i$, every site $y$, and every on-site operator $B_y$,
\begin{equation}
 \|[\tau_t(A),B_y]\|
 \le C\|A\|\|B_y\|e^{v|t|-c d_\Lambda(a_i,y)}.
 \label{eq:quasilocal-lr}
\end{equation}
The commutator is zero when $d_\Lambda(a_i,y)=\infty$.  Moreover,
\begin{equation}
 \|\tau_t(A)-E_{i,l}(\tau_t(A))\|
 \le \|A\|\min\{2,Ce^{v|t|-cl}\}.
 \label{eq:quasilocal-lr-ce}
\end{equation}
\end{lemma}

\begin{proof}
Fix $B_y$ and define
\[
 F(X,t)=\sup_{\substack{\supp A\subseteq X\\\|A\|\le1}}
                 \|[\tau_t(A),B_y]\|.
\]
For $t\ge0$, fix one such $A$, put $f(t)=[\tau_t(A),B_y]$, and let
$Q_X(t)=\sum_{j:X_j\cap X\ne\varnothing}\tau_t(h_j)$.
Since the terms disjoint from $X$ commute with $A$, differentiating and
using the Jacobi identity gives the exact equation
\[
 f'(t)=i[Q_X(t),f(t)]
       -i\sum_{j:X_j\cap X\ne\varnothing}
            [\tau_t(A),[\tau_t(h_j),B_y]].
\]
The propagator of the homogeneous equation is unitary conjugation.  Variation
of constants, followed by the supremum over $A$, consequently gives
\begin{equation}
 F(X,t)\le 2\|B_y\|\mathbf1_{\{y\in X\}}
   +2\sum_{j:X_j\cap X\ne\varnothing}\|h_j\|
                 \int_0^t F(X_j,u)\,du.
 \label{eq:quasilocal-lr-recursion}
\end{equation}
The same argument with $-H$ handles negative times.

Starting from $X=X_i$, iterate \eqref{eq:quasilocal-lr-recursion}.
At each step the sum of the available interaction norms is at most
$\kappa:=v_Rb_0$, by \eqref{eq:quasilocal-budget}.  A chain
$X_i,X_{j_1},\ldots,X_{j_m}$ with consecutive intersections and with
$y\in X_{j_m}$ has
$d_\Lambda(a_i,y)\le(m+1)R$.  The remainder after $M$ iterations is at
most $2\|B_y\|(2\kappa t)^M/M!$, since $F(X,u)\le2\|B_y\|$;
it tends to zero.  For every $\mu>0$, the resulting chain sum satisfies
\begin{align*}
 F(X_i,t)
 &\le 2\|B_y\|\sum_{m\ge0}
      \frac{(2\kappa t)^m}{m!}
      \mathbf1_{\{d_\Lambda(a_i,y)\le(m+1)R\}}\\
 &\le 2\|B_y\|e^{\mu R-\mu d_\Lambda(a_i,y)}
                    \exp(2\kappa e^{\mu R}t).
\end{align*}
This proves \eqref{eq:quasilocal-lr}.  If the distance is infinite, no
finite chain contributes and the vanishing remainder proves exact zero.

For completeness, write the expectation in \eqref{eq:quasilocal-ce}
as a product of on-site twirls $T_y$.  These are commuting contractions,
and telescoping that product shows, for any $B$,
\[
 \|B-E_{i,l}(B)\|
 \le\sum_{y\notin N_l(a_i)}\|B-T_y(B)\|
 \le\sum_{y\notin N_l(a_i)}\sup_{U_y}\|[B,U_y]\|,
\]
where $U_y$ ranges over on-site unitaries.  Sites in another component
contribute zero.  The number of sites at graph distance $d$ is at most
$C(d+1)^2$, because the graph ball lies in an ambient lattice ball.
Applying \eqref{eq:quasilocal-lr} and summing
$\sum_{d>l}(d+1)^2e^{-cd}\le C'e^{-c'l}$ proves the exponential bound
in \eqref{eq:quasilocal-lr-ce}.  Contractivity gives its other bound.
\end{proof}

The filter must remove ground-to-excited matrix elements exactly, while its
time kernel must decay fast enough to combine with
Lemma~\ref{lem:quasilocal-lr}.  The following elementary construction gives
every stretched exponential exponent $p/(p+1)$.

\begin{lemma}[A compactly supported spectral filter]
\label{lem:quasilocal-filter}
Fix an integer $p\ge1$, put $\alpha=p/(p+1)$, and let $\delta>0$.
Define the real even function
\[
 \chi(\omega)=
 \begin{cases}
  e\exp[-(1-(\omega/\delta)^2)^{-p}],&|\omega|<\delta,\\
  0,&|\omega|\ge\delta.
 \end{cases}
\]
Then $\chi\in C_c^\infty(\mathbb R)$, $\chi(0)=1$, and the real even
function
\[
 f(t)=\frac1{2\pi}\int_{\mathbb R}\chi(\omega)e^{-it\omega}\,d\omega
\]
satisfies
\begin{equation}
 \chi(\omega)=\int_{\mathbb R}f(t)e^{it\omega}\,dt,
 \qquad |f(t)|\le Ce^{-c|t|^\alpha}.
 \label{eq:quasilocal-filter-decay}
\end{equation}
Here $C,c>0$ depend only on $p,\delta$.
\end{lemma}

\begin{proof}
We give the derivative estimate responsible for the exponent.  For real
$u\in(-1,1)$ set $w=1-u^2$.  Choose $\eta_p>0$ sufficiently small.
On the complex disk $|z-u|\le\eta_pw$ one has
\[
 1-z^2=w(1+\zeta),\qquad |\zeta|\le3\eta_p,
 \qquad \operatorname{Re}(1+\zeta)^{-p}\ge c_p>0.
\]
The last inequality follows by continuity at $\zeta=0$, with $p$ fixed.
The function $e\exp[-(1-z^2)^{-p}]$ is analytic on the disk, and Cauchy's
estimate gives
\[
 \left|\frac{d^m}{du^m}\bigl(e\exp[-(1-u^2)^{-p}]\bigr)\right|
 \le C m!\eta_p^{-m}w^{-m}e^{-c_pw^{-p}}.
\]
For each fixed $m$ this tends to zero at both endpoints.  Extension by zero
is therefore smooth.  To make the estimate uniform in $u$, substitute
$s=w^{-p}$ and maximize $s^{m/p}e^{-c_ps}$.  Its maximum is bounded by
$C^{m+1}(m!)^{1/p}$, using $m!\ge(m/e)^m$; the case $m=0$ is absorbed
in $C$.  Scaling from $u$ to $\omega$ yields
\begin{equation}
 \|\chi^{(m)}\|_\infty\le C_1^{m+1}(m!)^{1+1/p}.
 \label{eq:quasilocal-gevrey}
\end{equation}

Integration by parts $m$ times has no boundary terms.  Consequently, for
$t\ne0$,
\[
 |f(t)|\le C\frac{C_1^m(m!)^{1+1/p}}{|t|^m}
       \le C\left(\frac{C_1m^{1+1/p}}{|t|}\right)^m.
\]
For sufficiently large $|t|$, choose
$m=\lfloor\eta |t|^{p/(p+1)}\rfloor$ with a fixed, sufficiently small
$\eta>0$.  The expression in parentheses is at most $1/2$, proving
the decay in \eqref{eq:quasilocal-filter-decay}; bounded $t$ is covered
by increasing $C$.  In particular $f\in L^1(\mathbb R)$.  Fourier
inversion gives the first identity in \eqref{eq:quasilocal-filter-decay}.
It can also be obtained directly by inserting $e^{-\varepsilon t^2}$,
interchanging the two integrals, and letting the resulting Gaussian
approximate identity converge to the continuous function $\chi$.
Reality and evenness follow from those properties of $\chi$.
\end{proof}

\subsection{Exact positive constraints and their local square roots}

We begin with Kitaev's centered spectral-filter construction of Hermitian
quasi-local constraints that annihilate the ground
state~\cite[Appendix~D.1.2, Proposition~D.1]{Kitaev2006Anyons}.
Its terms sum to the original Hamiltonian minus its ground energy.
Taking their absolute values gives the positive replacement used here:
the new sum dominates the original shifted Hamiltonian and therefore
retains a global gap.  The locality estimates below also quantify the
effect of this functional calculus.  The distinction between filtered
constraints and a positive sum with a gap is discussed by Swingle and
McGreevy~\cite[Section~8.4]{SwingleMcGreevy2016}.

We now combine the spectral and spatial estimates.  In the next statement,
the normalization $c_*$ and the gap lower bound $g$ may depend on the chosen
$p$.  Once $p$ is fixed, both are independent of the system size.

\begin{proposition}[Positive replacement of the Hamiltonian]
\label{prop:positive}
Suppose $H\Omega=E_0\Omega$, $\|\Omega\|=1$, and
\[
 H-E_0I\ge\Delta(I-P_\Omega),\qquad
 P_\Omega=|\Omega\rangle\langle\Omega|.
\]
For each integer $p\ge1$ and $\alpha=p/(p+1)$, there are operators
$k_i$, indexed by the original supports and retaining their anchors, and a
constant $c_*>0$ such that
\begin{equation}
 0\le k_i\le I,\qquad k_i\Omega=0,\qquad
 H_F:=\sum_{i\in\mathcal I}k_i
       \ge c_*^{-1}(H-E_0I)\ge g(I-P_\Omega),
 \qquad g:=\Delta/c_*.
 \label{eq:quasilocal-positive}
\end{equation}
Thus $H_F$ has the unique ground vector $\Omega$, ground energy zero, and
gap at least $g$.  Furthermore,
\begin{equation}
 k_{i,l}:=E_{i,l}(k_i)\in[0,I],\qquad
 \|k_i-k_{i,l}\|\le Ce^{-cl^\alpha}
 \quad(l\in\mathbb Z_{\ge0}).
 \label{eq:quasilocal-positive-tail}
\end{equation}
Each $k_i$ is supported in the connected component containing $a_i$.
\end{proposition}

\begin{proof}
Use Lemma~\ref{lem:quasilocal-filter} with $\delta=\Delta/2$, and define
\begin{equation}
 M_i=\int_{\mathbb R}f(t)\tau_t(h_i)\,dt
             -\langle\Omega,h_i\Omega\rangle I.
 \label{eq:quasilocal-centered-filter}
\end{equation}
The integral converges in norm and $M_i$ is Hermitian.  In an eigenbasis
of $H$, its uncentered ground-to-energy-$E$ matrix element equals the
corresponding matrix element of $h_i$ multiplied by $\chi(E-E_0)$.
All excited energies satisfy $E-E_0\ge\Delta$, where $\chi$ vanishes.
The ground component is unchanged because $\chi(0)=1$, and the scalar
in \eqref{eq:quasilocal-centered-filter} removes it.  Hence $M_i\Omega=0$.
Also, since $\tau_t(H)=H$ and $\int f=1$,
\begin{equation}
 \sum_i M_i=H-E_0I,\qquad
 \|M_i\|\le J(\|f\|_1+1).
 \label{eq:quasilocal-filter-sum}
\end{equation}
Choose $c_*:=\max\{1,J(\|f\|_1+1)\}$ and set $k_i=|M_i|/c_*$.
For a Hermitian operator, $|M_i|\ge M_i$ and
$\ker|M_i|=\ker M_i$.  These observations and
\eqref{eq:quasilocal-filter-sum} prove \eqref{eq:quasilocal-positive}.

It remains to check locality after taking the absolute value.  The scalar
in \eqref{eq:quasilocal-centered-filter} is fixed by $E_{i,l}$.  For
$|t|\le\varepsilon l$, choose $\varepsilon>0$ small enough that
\eqref{eq:quasilocal-lr-ce} is at most $Ce^{-c'l}$ for $A=h_i$.
For the remaining times use its bound $2J$.  Thus
\begin{align}
 \|M_i-E_{i,l}(M_i)\|
 &\le Ce^{-c'l}\|f\|_1
       +2J\int_{|t|>\varepsilon l}|f(t)|\,dt
   \le C'e^{-c''l^\alpha}.
 \label{eq:quasilocal-filter-locality}
\end{align}
The last tail estimate follows by substituting $u=|t|^\alpha$ and
absorbing the resulting fixed power of $u$ into the exponential.  The
case $l=0$ follows from the uniform norm bound.

We use the dimension-independent inequality
\begin{equation}
 \|U^{1/2}-V^{1/2}\|\le\|U-V\|^{1/2}
 \qquad(U,V\ge0).
 \label{eq:quasilocal-square-root}
\end{equation}
Here is an order proof.  The integral formula
\[
 T^{1/2}=\frac1\pi\int_0^\infty
                  T(T+sI)^{-1}s^{-1/2}\,ds\qquad(T\ge0)
\]
follows from scalar integration and the spectral theorem.  It shows that
the square root is operator monotone: inversion reverses the order of
positive invertible operators, and
$T(T+sI)^{-1}=I-s(T+sI)^{-1}$.  With
$\epsilon=\|U-V\|$, monotonicity and scalar functional calculus give
$U^{1/2}\le(V+\epsilon I)^{1/2}\le V^{1/2}+\sqrt\epsilon I$.
Interchanging $U,V$ proves \eqref{eq:quasilocal-square-root}.

Apply this with $U=M_i^2$ and $V=(E_{i,l}M_i)^2$.  Both factors in the
latter square are Hermitian, and their norms are at most $c_*$.  Therefore
\[
 \big\||M_i|-|E_{i,l}M_i|\big\|
 \le\bigl(2c_*\|M_i-E_{i,l}M_i\|\bigr)^{1/2}
 \le Ce^{-cl^\alpha}.
\]
The operator $|E_{i,l}M_i|$ is supported in $N_l(a_i)$.  Since $E_{i,l}$
fixes it and is contractive,
\[
 \||M_i|-E_{i,l}(|M_i|)\|
       \le2\big\||M_i|-|E_{i,l}M_i|\big\|.
\]
After division by $c_*$ this proves
\eqref{eq:quasilocal-positive-tail}; its order bounds follow from
positivity and unitality of $E_{i,l}$.

Finally, every original support lies in one component.  The Hamiltonian
is a tensor sum over components, so $\tau_t(h_i)$ is supported in the
component of $a_i$.  The integral, scalar centering, and functional calculus
in the above construction all stay in that component's operator algebra.
This proves the exact component assertion.
\end{proof}

The same construction supplies local approximants to the square roots of
the positive constraints.  We include the associated channels because
their unitality will allow spatial localization of products in
Section~\ref{sec:amplification}.

\begin{lemma}[Square roots and local channels]
\label{lem:quasilocal-roots}
For the constraints of Proposition~\ref{prop:positive}, define
\[
 G_i=(I-k_i)^{1/2},\quad K_i=k_i^{1/2},\qquad
 G_{i,l}=(I-k_{i,l})^{1/2},\quad K_{i,l}=k_{i,l}^{1/2}.
\]
These are positive contractions, $G_i\Omega=\Omega$, $K_i\Omega=0$,
and
\begin{equation}
 \|G_i-G_{i,l}\|+\|K_i-K_{i,l}\|\le Ce^{-cl^\alpha},
 \qquad G_{i,l}^2+K_{i,l}^2=I.
 \label{eq:quasilocal-root-tail}
\end{equation}
Define maps on the full operator algebra by
\[
 \mathcal E_i(B)=G_iBG_i+K_iBK_i,\qquad
 \mathcal E_{i,l}(B)=G_{i,l}BG_{i,l}+K_{i,l}BK_{i,l}.
\]
Both maps are unital and completely positive, and are contractions in
operator norm, also after tensoring with the identity map on any auxiliary
matrix algebra.  The completely bounded norm, denoted by
$\|\cdot\|_{\mathrm{cb}}$, is the supremum of the operator norms on all
such auxiliary extensions.  In this norm,
\begin{equation}
 \|\mathcal E_i-\mathcal E_{i,l}\|_{\mathrm{cb}}
          \le Ce^{-cl^\alpha}.
 \label{eq:quasilocal-channel-tail}
\end{equation}
The map $\mathcal E_{i,l}$ fixes every operator commuting with the full
algebra on $N_l(a_i)$, including operators supported outside that ball.
All these conclusions allow arbitrary spectator auxiliary systems.
\end{lemma}

\begin{proof}
Apply \eqref{eq:quasilocal-square-root} to $k_i,k_{i,l}$ and to
$I-k_i,I-k_{i,l}$, using \eqref{eq:quasilocal-positive-tail}.
The order and ground-vector assertions follow from functional calculus.
Each displayed channel is a sum of two maps of the form $B\mapsto A^*BA$;
this proves complete positivity, while the sum of the squared Kraus
operators is $I$.  More explicitly, the map
$v\mapsto G_iv\oplus K_iv$ is an isometry and
$\mathcal E_i(B)$ is the compression of $B\oplus B$ by that isometry.
This proves contractivity, and the same proof works for the local maps and
after adjoining an identity tensor factor.  Finally,
\[
 \|ABA-A'BA'\|
 \le(\|A\|+\|A'\|)\|A-A'\|\|B\|,
\]
also on every auxiliary extension.  Applying this to both summands proves
\eqref{eq:quasilocal-channel-tail}, where the completely bounded norm is
the supremum of these extended operator norms.  An operator commuting
with the ball algebra commutes with both local square roots, so its image
is itself times $G_{i,l}^2+K_{i,l}^2=I$.
\end{proof}

\subsection{Truncation with an error independent of the total volume}

We next make the positive constraints strictly local near a testing set.
Conditional expectation preserves positivity but need not preserve
$k_i\Omega=0$.  We therefore bound the norm of the total perturbation and
work with the truncated Hamiltonian's exact ground vector.  Farther
constraints receive larger truncation radii, so that summing all errors
does not charge the entire volume.  Graph balls will also let us count the
terms crossing a cut by
the number of edges of that cut.  The choice $p=1$ suffices here: tails
of the form $e^{-c\sqrt l}$ give radii of order $(\log n)^2$ for
inverse-power accuracy at scale $n$.

\begin{proposition}[Truncation near a finite set]
\label{prop:truncation}
Use Proposition~\ref{prop:positive} with $p=1$, so $\alpha=1/2$.
Fix $C_0>0$.  There is a constant $C_1$, depending only on
$q,R,J,\Delta,C_0$, with the following property for every real $n\ge2$
and every nonempty $S_0\subseteq\Lambda$ with $|S_0|\le C_0n^2$.
Put $r_0=\lceil C_1(\log n)^2\rceil$.  If
$d_i^0:=d_\Lambda(a_i,S_0)<\infty$, define
\begin{equation}
 r_i=\max\{r_0,\lfloor d_i^0/2\rfloor\},\qquad
 \widetilde h_i=k_{i,r_i},\qquad \widetilde X_i=N_{r_i}(a_i).
 \label{eq:quasilocal-variable-radius}
\end{equation}
If $d_i^0=\infty$, set $\widetilde h_i=k_i$ and designate its anchor's
entire component as $\widetilde X_i$.  Then
$\widetilde H=\sum_i\widetilde h_i$ is a sum of positive contractions and
\begin{equation}
 \|\widetilde H-H_F\|\le\epsilon_n,
 \qquad \epsilon_n:=\min\{n^{-1000},g/4\}.
 \label{eq:quasilocal-truncation-error}
\end{equation}
Its ground energy $\widetilde E_0$ lies in $[0,\epsilon_n]$, its ground
vector $\widetilde\Omega$ is unique, and its gap is at least $g/2$.
After a choice of phase,
\begin{equation}
 \|\widetilde\Omega-\Omega\|\le2\sqrt{\epsilon_n/g},\qquad
 \frac12\big\||\widetilde\Omega\rangle\langle\widetilde\Omega|
                 -P_\Omega\big\|_1\le\sqrt{2\epsilon_n/g}.
 \label{eq:quasilocal-ground-distance}
\end{equation}

For a set $X\subseteq S_0$, let $b_X=|\partial_\Lambda X|$ and let
$\mathcal C_X$ be the labels whose designated supports
$\widetilde X_i$ meet both $X$ and $\Lambda\setminus X$.
Every such support has radius exactly $r_0$, and
\begin{equation}
 |\mathcal C_X|\le Cb_X(r_0+1)^2,\qquad
 \log d_i\le C(r_0+1)^2\quad(i\in\mathcal C_X),\qquad
 d_i:=q^{|\widetilde X_i|}.
 \label{eq:quasilocal-crossing-count}
\end{equation}
Consequently the cut parameter
\begin{equation}
 \mathcal B_X:=1+\sum_{i\in\mathcal C_X}\log^2(e d_i)
       \le C(1+b_X)(r_0+1)^6.
 \label{eq:quasilocal-cut-budget}
\end{equation}
In particular, $b_X\le C_2nD$ and $D\ge1$ imply
$\mathcal B_X\le CnD(\log n)^{12}$, with the constant also depending
on $C_2$.
\end{proposition}

\begin{proof}
We first prove the summation estimate for any fixed $\alpha\in(0,1)$
for which \eqref{eq:quasilocal-positive-tail} holds, and any integer
$r_0\ge1$.  At each finite distance $d$ from
$S_0$ there are at most $C\mu_R|S_0|(d+1)^2$ labels: every anchor at that
distance belongs to $N_d(x)$ for some $x\in S_0$, and anchor multiplicity
is bounded by $\mu_R$.  Labels at infinite distance have no error by
definition.  The triangle inequality therefore gives
\begin{align}
 \|\widetilde H-H_F\|
 &\le C|S_0|\sum_{d\ge0}(d+1)^2
       \exp\bigl(-c\max\{r_0,\lfloor d/2\rfloor\}^{\alpha}\bigr)
   \le C'|S_0|e^{-c'r_0^\alpha}.
 \label{eq:quasilocal-volume-free}
\end{align}
To justify the final inequality, the terms $d\le2r_0+2$ have sum at
most $C(r_0+1)^3e^{-cr_0^\alpha}$.  For $d>2r_0+2$ use
$\lfloor d/2\rfloor\ge d/3$; comparison with the integral of
$(x+1)^2e^{-c(x/3)^\alpha}$, followed by $u=x^\alpha$, bounds the
remaining tail by a fixed power of $r_0+1$ times $e^{-c''r_0^\alpha}$.
All such powers are absorbed by decreasing the positive exponential
constant.  Thus \eqref{eq:quasilocal-volume-free} has no factor involving
$|\Lambda|$.

Now take $\alpha=1/2$, $|S_0|\le C_0n^2$, and
$r_0=\lceil C_1(\log n)^2\rceil$.  The right-hand side of
\eqref{eq:quasilocal-volume-free} is at most
$C'C_0n^{2-c'\sqrt{C_1}}$.  Taking $C_1$ sufficiently large proves
\eqref{eq:quasilocal-truncation-error}, including its $g/4$ bound for
all $n\ge2$.

Positivity gives $\widetilde E_0\ge0$, and testing
$\widetilde H$ on $\Omega$ gives $\widetilde E_0\le\epsilon_n$.
The min--max principle and \eqref{eq:quasilocal-truncation-error} put
the second eigenvalue of $\widetilde H$ at least $g-\epsilon_n$.
The difference between its first two eigenvalues is consequently at least
$g-2\epsilon_n\ge g/2$, which also proves uniqueness.  Moreover,
\[
 g\bigl(1-|\langle\Omega,\widetilde\Omega\rangle|^2\bigr)
 \le\langle\widetilde\Omega,H_F\widetilde\Omega\rangle
 \le\widetilde E_0+\epsilon_n\le2\epsilon_n.
\]
The trace distance of two pure states equals the square root of one minus
their squared overlap, giving the second estimate in
\eqref{eq:quasilocal-ground-distance}.  Choose the overlap real and
nonnegative, and use $1-s\le1-s^2$ for $0\le s\le1$ to obtain the
first estimate.

Finally, if a designated support meets $X\subseteq S_0$, then its anchor
has finite distance from $S_0$.  If $r_i>r_0$, formula
\eqref{eq:quasilocal-variable-radius} gives
$r_i=\lfloor d_i^0/2\rfloor<d_i^0$, so that its ball cannot meet
$S_0$.  Hence every label in $\mathcal C_X$ has $r_i=r_0$.
If its anchor lies in $X$, take a shortest path from the anchor to a
point of the ball outside $X$; otherwise take one to a point inside $X$.
This path stays in the ball and contains an edge of $\partial_\Lambda X$.
Thus the anchor is within $r_0$ of an endpoint of a crossing edge.
There are at most $2b_X$ endpoints, at most $C(r_0+1)^2$ sites in each
such ball, and at most $\mu_R$ labels per anchor.  This proves the first
bound in \eqref{eq:quasilocal-crossing-count}.  The ball volume bound
also gives $\log d_i=|\widetilde X_i|\log q\le C(r_0+1)^2$.
Summing its squared version proves \eqref{eq:quasilocal-cut-budget}.
Since $r_0+1\le C(\log n)^2$ for $n\ge2$, its stated consequence
follows.
\end{proof}

In the geometric applications of Section~\ref{sec:scanner}, a finite
$T\subset\mathbb Z^2$ and an integer $L\ge0$ will be specified.  With
$T_L=T+([-L,L]^2\cap\mathbb Z^2)$, we take $S_0=A\cap T_L$.
Only the bound $|S_0|\le Cn^2$ is needed in
Proposition~\ref{prop:truncation}; the geometric estimates will provide
$b_X=O(nD)$ for the cuts that arise there.
If $S_0$ is empty we retain all constraints exactly, and every cut
$X\subseteq S_0$ has no crossing label.  More generally, components that
do not meet $S_0$ are retained exactly throughout the construction.
The exponent $1000$ in the approximation error can be replaced by any
fixed positive exponent by increasing $C_1$; the powers of $\log n$ in
the support and cut bounds do not change.
